As expansões hiperbólicas são obtidas analogamente, dividindo as séries de $\sinh x$ e $\cosh x$:
\[
\sinh x = x + \frac{x^{3}}{6} + \frac{x^{5}}{120} + \cdots, \qquad \cosh x = 1 + \frac{x^{2}}{2} + \frac{x^{4}}{24} + \cdots.
\]
O contraste com as séries trigonométricas é que os sinais são todos positivos (sem alternância). Dividindo:
\[
\tanh x = x - \frac{x^{3}}{3} + \frac{2x^{5}}{15} - \frac{17x^{7}}{315} + \cdots,
\]
\[
\coth x = \frac{1}{x} + \frac{x}{3} - \frac{x^{3}}{45} + \frac{2x^{5}}{945} + \cdots.
\]
Pode-se verificar a relação com as expansões de $\tan$ e $\cot$: substituição $x\to ix$ transforma uma nas outras, usando $\tanh(ix)=i\tan x$.